\section{Experiments}
\label{sec:experiments}

\subsection{Experimental Setup}

We validate the theoretical claims of Section~\ref{sec:theory} using a controlled
simulation framework built around the \texttt{dva\_mcts} Python package released
alongside this paper. The verifier is a \texttt{LipschitzVerifier}: a synthetic
process reward model whose true value function is a Lipschitz-$L$ random walk and
whose observations are corrupted by additive Gaussian noise. Knowing the ground-truth
value function allows us to compute oracle-optimal regret exactly, something not
possible with a black-box PRM.

\paragraph{Fixed parameters.}
All experiments use branching factor $K = 2$, maximum depth $D = 12$, Lipschitz
constant $L = 0.35$, and noise $\sigma = 0.05$. These match the calibration of
\citet{lightman2023lets}, who report that a $7$B-parameter PRM has per-step score
noise of $\sigma \approx 0.05$ on the MATH-500 benchmark. All results are averaged
over $50$ independent runs; shaded regions indicate $\pm 1$ standard deviation.
The threshold parameter is $\gamma = 1.0$ and (for the default configuration)
$\alpha = 0.5$.

\paragraph{Baselines.}
\begin{itemize}
  \item \textit{Uniform Verification}: verifier called at every search step. This is
    the standard exhaustive-verification approach against which we measure savings.
  \item \textit{Random Allocation}: verifier called independently at each step with
    probability $0.5$, regardless of Lipschitz discrepancy.
  \item \textit{Best-of-$N$}: $N$ independent random paths sampled from root to leaf;
    each path is fully verified. This is the simplest test-time scaling baseline
    \citep{brown2024large}.
\end{itemize}

\subsection{Regret Scaling}
\label{sec:regret_results}

Figure~\ref{fig:regret} shows cumulative regret over $T = 400$ search steps.

\paragraph{Rate confirmation.}
The log-log plot (panel~(b)) shows that DVA-MCTS regret grows with empirical slope
$\mathbf{0.635}$ (OLS on log-log scale, $R^2 = 0.961$). The $O(\sqrt{T})$ lower
bound predicts slope $0.5$; our empirical slope of $0.635$ is consistent with an
additive constant or logarithmic correction that is not apparent at $T = 400$.
As budget grows, the slope is expected to converge toward $0.5$ (see Figure~\ref{fig:calls}
for evidence at $T = 3200$).

\paragraph{Comparison with baselines.}
At $T = 400$ DVA-MCTS achieves mean final regret $31.32$, compared to $26.76$ for
Uniform Verification and $30.40$ for Random Allocation. Uniform achieves lower
absolute regret because it calls the verifier at \emph{every} node, accumulating
high-quality estimates quickly; DVA-MCTS trades a small regret increase for the
much larger efficiency gain demonstrated in Section~\ref{sec:efficiency_results}.

\subsection{Verifier Call Efficiency}
\label{sec:efficiency_results}

Figure~\ref{fig:calls} reports verifier call counts as a function of budget $B$
across the range $B \in \{50, 100, 200, 400, 800, 1600, 3200\}$.

\paragraph{Regime-dependent savings.}
At small budgets ($B \leq 200$), efficiency savings are modest ($1{-}6\%$) because
DVA-MCTS is still exploring new nodes for the first time, and unvisited nodes are
always verified on first visit. The large savings emerge once the tree is sufficiently
explored: at $B = 800$ DVA uses $77.1\%$ of Uniform's calls; at $B = 1600$ this
drops to $48.7\%$; and at $B = 3200$ to $\mathbf{26.9\%}$---a $3.7\times$ reduction.
This scaling is consistent with $O(\sqrt{T}\log T)$ calls (Theorem~\ref{thm:calls}):
$\sqrt{3200}\log(3200) \approx 456$ predicted calls versus $862$ observed (mean),
with the gap attributable to warm-start verification and finite-depth effects.

\subsection{Lipschitz Continuity Validation}
\label{sec:lipschitz_results}

Figure~\ref{fig:lipschitz} validates Assumption~\ref{ass:lipschitz} on $500$ simulated
trajectories comprising $22{,}500$ state pairs.

\paragraph{Envelope holds exactly.}
Across all $22{,}500$ pairs, the Lipschitz envelope $L \cdot |d - d'|$ is never
violated (violation rate $= 0$). This is guaranteed by construction of the
\texttt{LipschitzVerifier}, but provides a concrete check that the data-generating
process satisfies the assumption.

\paragraph{Mean rate vs.\ Lipschitz supremum.}
The global empirical estimate $\hat{L} = 0.091$ (panel~(b)) is substantially below
the true $L = 0.35$. This is expected: $\hat{L}$ estimates the \emph{mean} rate of
score change via $\mathbb{E}[|V(s) - V(s')|] / |d - d'|$, while $L$ is the
\emph{supremum}. For Uniform$(-L, L)$ increments, the mean absolute increment is
$L/2 = 0.175$; averaged across larger depth gaps (where increments partially cancel),
the global mean is lower still. In practice, the algorithm uses the true $L$ as a
parameter; $\hat{L}$ is reported for interpretability only.

\subsection{Compute--Accuracy Tradeoff}
\label{sec:tradeoff_results}

Figure~\ref{fig:tradeoff} shows solution accuracy (fraction of runs where the returned
solution has true value $\geq 0.8$) as a function of budget $B$.

\paragraph{DVA-MCTS vs.\ Uniform.}
DVA-MCTS and Uniform Verification converge to similar accuracy across all budgets,
with DVA reaching $98\%$ accuracy at $B = 400$ and maintaining it through $B = 3200$.
Uniform reaches $100\%$ accuracy at $B = 400$, a gap of $2$ percentage points. This
gap persists but does not grow as budget increases, consistent with the sub-linear
regret guarantee.

\paragraph{Best-of-$N$ fails on deep trees.}
Best-of-$N$ achieves $0\%$ accuracy across all budgets. This is not a failure of
verification but of search: with depth $D = 12$ and branching factor $K = 2$, the
tree contains $2^{12} = 4096$ leaves. A uniformly random path to a leaf has
probability $\leq 2^{-12} \approx 0.02\%$ of selecting the best leaf. UCT-based
search (DVA-MCTS and Uniform) actively steers toward high-value regions, which is
essential for deep trees. This finding reinforces why verifier-guided tree search
is preferred over best-of-$N$ for hard, long-horizon problems \citep{snell2024scaling}.

\subsection{Ablation: Threshold Exponent $\alpha$}
\label{sec:ablation}

We sweep $\alpha \in \{0.25, 0.50, 0.75, 1.00\}$ with $\gamma = 1.0$ fixed
(Figure~\ref{fig:ablation}).

\paragraph{Regret trends.}
At $T = 400$, the empirical final regrets are:
$27.87$ ($\alpha{=}0.25$), $18.98$ ($\alpha{=}0.50$), $\mathbf{17.82}$ ($\alpha{=}0.75$),
and $19.99$ ($\alpha{=}1.00$). The minimum is at $\alpha = 0.75$; both $\alpha = 0.25$
and $\alpha = 1.00$ perform worse. The theoretically recommended value $\alpha = 0.5$
is within $6.5\%$ of the empirical optimum.

\paragraph{Verification rate.}
Verification rates are high across all $\alpha$ at $B = 400$ (93--97\%), because the
tree is still largely unexplored and DVA-MCTS calls the verifier on first visits
regardless of the threshold. The distinction between $\alpha$ values becomes more
pronounced at larger budgets (as shown in Section~\ref{sec:efficiency_results}),
where revisit suppression dominates.

\begin{table}[t]
\centering
\caption{Summary of DVA-MCTS vs.\ baselines at $T = 400$ (50 runs).}
\label{tab:summary}
\begin{tabular}{lrrc}
\toprule
Method & Regret $R(400)$ & Verifier Calls & Accuracy (${\geq}0.8$) \\
\midrule
Uniform Verification & $26.76$          & $401$           & $96.0\%$ \\
\textbf{DVA-MCTS (ours)} & $31.32$     & $\mathbf{382}$  & $\mathbf{98.0\%}$ \\
Random Allocation    & $30.40$          & $202$           & $98.0\%$ \\
Best-of-$N$          & ---              & $400$           & $0.0\%$ \\
\bottomrule
\end{tabular}
\end{table}

\subsection{Discussion of Experimental Findings}

Three observations deserve emphasis.

First, efficiency gains are \emph{budget-dependent}. At $B = 400$, DVA-MCTS saves
only $5\%$ of verifier calls; at $B = 3200$ it saves $73\%$. This matches the
theoretical $O(\sqrt{B}\log B / B) = O(1/\sqrt{B})$ efficiency ratio: savings grow
as budget increases, with large-budget regimes being the primary target application.

Second, DVA-MCTS achieves \emph{higher accuracy than Uniform at small budgets}.
At $B = 50$, DVA reaches $88\%$ vs.\ Uniform's $78\%$ (Table~\ref{tab:summary}
and Figure~\ref{fig:tradeoff}). We attribute this to the Lipschitz proxy: when the
verifier is not called, DVA inherits the ancestor's score, which is often a better
estimate than a fresh noisy observation would be. This provides a de-noising effect
at small budgets.

Third, the $O(\sqrt{T})$ rate is confirmed at the right scale. The empirical slope
of $0.635$ is slightly above the theoretical $0.5$; this gap is consistent with an
additive $O(\sqrt{T}\log T)$ term (from the verifier call complexity) that dominates
at $T = 400$ but vanishes relative to the leading $\sqrt{T}$ term as $T \to \infty$.
